% ═══════════════════════════════════════════════════════════════════
% §4  정상상태 급탕 시스템 모델
% Source: enex_engine.py
% ═══════════════════════════════════════════════════════════════════
\section{정상상태 급탕 시스템 모델}
\label{sec:enex-engine}

본 절은 \texttt{enex\_engine.py}의 정상상태 급탕 시스템 모델에 대한 에너지, 엔트로피, 엑서지 수지를 제시한다. 각 모델은 \emph{온수 저탕조}와 \emph{혼합 밸브}로 구성되며, 열원의 종류에 따라 모델이 구분된다.

\subsection{공통 시스템 경계}
\label{sec:dhw-common}

모든 급탕 모델은 두 구성 요소의 경계를 공유한다:
\begin{enumerate}
  \item \textbf{온수 저탕조}: 열을 입력받아 $T_\text{tank}$의 물을 저장하고, 주위로 열손실이 발생한다.
  \item \textbf{혼합 밸브}: 고온의 저탕조 물과 저온의 급수를 혼합하여 급탕 온도 $T_\text{serv}$로 공급한다.
\end{enumerate}

혼합비 $\alpha$는 저탕조 취수 비율을 결정한다:
\begin{equation}\label{eq:dhw-alpha}
  \alpha = \frac{T_\text{serv} - T_\text{sup}}{T_\text{tank} - T_\text{sup}}
  \,,\quad
  \dot{V}_\text{sup,tank} = \alpha\,\dot{V}_\text{serv}
  \,,\quad
  \dot{V}_\text{sup,mix} = (1 - \alpha)\,\dot{V}_\text{serv}
\end{equation}

\subsection{공통 열역학 정식화}

\paragraph{물 유동의 에너지.}
\begin{equation}\label{eq:dhw-Q}
  Q_w = \rho_w c_w \dot{V}\,(T - T_0)
\end{equation}

\paragraph{물 유동의 엔트로피.}
\begin{equation}\label{eq:dhw-S}
  \dot{S}_w = \rho_w c_w \dot{V}\,\ln\frac{T}{T_0}
\end{equation}

\paragraph{물 유동의 엑서지 (비압축성 유체).}
\begin{equation}\label{eq:dhw-X}
  X_w = \rho_w c_w \dot{V}\left[(T - T_0) - T_0\,\ln\frac{T}{T_0}\right]
\end{equation}

\paragraph{저탕조 열손실.}
\begin{equation}\label{eq:dhw-Ql}
  Q_\text{l,tank} = U_\text{tank}\,(T_\text{tank} - T_0)
\end{equation}
여기서 $U_\text{tank}$ [W/K]은 원통형 외피 및 단열재의 열저항으로부터 산출된다.

\paragraph{엔트로피 생성 및 엑서지 파괴.}
각 구성 요소 $k$에 대해:
\begin{equation}\label{eq:dhw-Sg}
  \dot{S}_{\text{g},k} = \sum \dot{S}_\text{out} - \sum \dot{S}_\text{in}
  \,,\qquad
  X_{\text{c},k} = \dot{S}_{\text{g},k}\,T_0
\end{equation}

\subsection{전기 보일러}

열 입력은 전기 소비 전력과 동일하다:
\begin{equation}\label{eq:eb-energy}
  E_\text{heater} = Q_\text{w,tank} + Q_\text{l,tank} - Q_\text{w,sup,tank}
\end{equation}
전기는 순수 일($T_\text{source} \to \infty$)이므로 $\dot{S}_\text{heater} = 0$이고 $X_\text{heater} = E_\text{heater}$이다.

\subsection{가스 보일러}

가스 입력은 연소 비가역성을 수반한다:
\begin{align}
  Q_\text{gas} &= \frac{Q_\text{w,tank} + Q_\text{l,tank} - Q_\text{w,sup,tank}}{\eta_\text{gas}} \label{eq:gb-Qgas} \\
  \dot{S}_\text{gas} &= \frac{Q_\text{gas}}{T_\text{ad}} \label{eq:gb-Sgas} \\
  X_\text{gas} &= Q_\text{gas}\left(1 - \frac{T_0}{T_\text{ad}}\right) \label{eq:gb-Xgas}
\end{align}
여기서 $\eta_\text{gas}$는 버너 효율, $T_\text{ad}$는 천연가스의 단열 화염 온도이다.

\subsection{히트펌프 급탕기}

히트펌프 모델은 ASHP COP 상관식을 사용한다 (\ref{sec:enex-functions}절 참조):
\begin{align}
  \text{COP} &= f(T_0,\,T_\text{tank}) \label{eq:hpb-COP}\\
  E_\text{hp} &= \frac{Q_\text{w,tank} + Q_\text{l,tank} - Q_\text{w,sup,tank}}{\text{COP}} \label{eq:hpb-Ehp}\\
  Q_\text{evap} &= Q_\text{cond} - E_\text{hp} \label{eq:hpb-Qevap}
\end{align}
팬 및 압축기 전력은 일 입력($\dot{S} = 0$)으로 모델링된다.

\subsection{엑서지 효율 (전 모델 공통)}

\begin{equation}\label{eq:dhw-Xeff}
  \eta_X = \frac{X_\text{w,serv}}{X_\text{input}}
\end{equation}
여기서 $X_\text{input}$은 모델 유형에 따라 $X_\text{heater}$, $X_\text{gas}$, 또는 $X_\text{hp} + X_\text{fan}$이다.
